\section{Introduction}
\label{sec:introduction}

Modern networks are increasingly expected to support rapid configuration changes, stringent service requirements, and automated control. Testing these changes directly on operational infrastructure is costly and can be unsafe. A network digital twin (NDT) addresses this problem by providing a digital representation in which network state, models, and control interfaces can be used to evaluate behaviour before or alongside physical deployment. Current reference work describes NDTs through data, models, interfaces, and mappings, and presents closed-loop management as a central objective \citep{zhou2026ndt,itut2022y3090}. Data-driven NDT research further emphasizes fidelity, efficiency, and flexibility in network performance modelling \citep{almasan2022ndt,hui2023digitaltwin}.

The practical challenge is not simply to instantiate a topology. A useful experimental twin must transform a scenario into a controlled execution, verify routing and impairment configuration, collect measurements, and retain enough provenance to support later analysis. Lightweight emulation established that real network applications and protocol stacks can be exercised on a single host \citep{lantz2010network}, while container-based emulation showed how runnable specifications and retained artifacts can improve reproducibility \citep{handigol2012reproducible}. However, architecture, emulation, topology scaling, AI-assisted configuration, and learning-based control are commonly evaluated as separate concerns. This fragmentation makes it difficult to assess whether they can share one bounded execution and evidence workflow without overstating fidelity or automation.

This dissertation addresses that integration problem through a scenario-driven Docker NDT. JSON scenarios are validated before execution; Docker containers and networks implement endpoint and router roles; deterministic forwarding and routes establish reachability; Linux \texttt{tc}, \texttt{netem}, and token-bucket filtering impose bandwidth, delay, and loss; and ping and iperf3 produce measurements that are retained as structured evidence. The same guarded execution environment is extended to a realistic Germany50-derived topology, validated AI-assisted scenario generation, a real-Docker control loop, and a deliberately bounded interactive Dashboard.

\subsection{Research gap and questions}

Existing work separately provides NDT architectures, reproducible/container-based emulation, realistic topology datasets, AI-assisted networking, and digital-twin-assisted reinforcement learning. The research gap considered here is the integrated, quantitative evaluation of these elements in a common scenario-driven Docker workflow. The work does not claim to be the first network emulator, NDT, AI networking system, or NDT/RL system. Instead, it asks whether a compact implementation can enforce controlled conditions, scale its topology and routing mechanisms, support guarded automation, and preserve claim-ready evidence.

The main research question is:

\begin{quote}
\emph{How effectively can a Docker-based network digital twin enforce and measure controlled network conditions while remaining reproducible and extensible across increasingly complex topologies and automated control workflows?}
\end{quote}

It is decomposed into three questions:

\begin{description}
\item[RQ1:] How closely and consistently do measured network behaviours follow configured bandwidth, latency, and packet-loss conditions across repeated experiments?
\item[RQ2:] How does the framework extend from simple Docker topologies to multi-router and larger network topologies while preserving validated routing and measurable real traffic?
\item[RQ3:] Can validated AI-assisted scenario generation and closed-loop learning-based control operate on the same digital-twin execution environment, and what limitations emerge in doing so?
\end{description}

\subsection{Contributions}

The dissertation makes three bounded contributions. First, it presents a reproducible, scenario-driven Docker NDT framework that joins validation, topology construction, routing, impairment, measurement, cleanup, and evidence retention. Second, it provides a quantitative evaluation of controlled bandwidth, one-way delay, packet loss, multi-router execution, and Germany50 topology scaling. Third, it evaluates guarded AI-assisted scenario generation and a real-Docker closed-loop controller, including the negative result that the evaluated threshold heuristic outperformed Q-learning. These contributions are limited to the retained configurations: Germany50 real traffic covers three representative paths, the full route plan is dry-run evidence, provider outcomes remain separate, and the Dashboard is a supporting small-topology interface.

The remainder of this dissertation reviews related work, describes the system and experimental methods, reports the retained results, discusses the research questions and threats to validity, and concludes with the scope of the supported claims.
